\par\Needspace{6\baselineskip}\medskip
\begin{center}\hypertarget{native-234-article169}{}\pdfbookmark[2]{169.}{article-169}169.\end{center}
Nach Abschluss der Rechnung des vorigen Artikels ist zunächst der Bogen $C'C''$ zu bestimmen. Am bequemsten geschieht dies, wenn man zuvor wie in Artikel 137 den Schnittpunkt $D$ der Großkreise $A'C'B'$, $A''C''B''$ und ihre gegenseitige Neigung $\epsilon$ ermittelt. Dann erhält man aus $\epsilon$, $C'D=z'+B'D$ und $C''D=z''+B''D$ mit denselben Formeln, die wir in Artikel 144 angegeben haben, sowohl $C'C''=v''-v'$ als auch die Winkel $(u',u'')$, unter denen die Großkreise $A'B'$, $A''B''$ den Großkreis $C'C''$ schneiden.

Sobald der Bogen $v''-v'$ gefunden ist, ergeben sich $v'-v$ und $r$ aus der Verbindung der Gleichungen
\[
r\sin(v'-v)=\frac{r''\sin(v''-v')}{P'}
\]
\[
r\sin(v'-v+v''-v')=\frac{1+P'}{P'}\cdot\frac{r'\sin(v''-v')}{1+\dfrac{Q'}{r'^3}}
\]
und ebenso $r'''$ und $v'''-v''$ aus der Verbindung der folgenden Gleichungen
\[
r'''\sin(v'''-v'')=\frac{r'\sin(v''-v')}{P''}
\]
\[
r'''\sin(v'''-v''+v''-v')=\frac{1+P''}{P''}\cdot\frac{r''\sin(v''-v')}{1+\dfrac{Q''}{r''^3}}\mathpunct{\cdot}
\]
\SourceLeaf{235}Alle auf diese Weise gefundenen Zahlen wären exakt, hätte man von Anfang an von den wahren Werten von $P',P'',Q',Q''$ ausgehen können. Dann wäre wie in Artikel 149 die Lage der Bahnebene entweder aus $A'C',u'$ und $\gamma'$ oder aus $A''C'',u''$ und $\gamma''$ zu bestimmen, und die Abmessungen der Bahn selbst entweder aus $r',r'',t',t''$ und $v''-v'$ oder, was genauer ist, aus $r,r''',t,t'''$ und $v'''-v$. Bei der ersten Rechnung übergehen wir jedoch all dies und bemühen uns vor allem um bessere Näherungswerte für die Größen $P',P'',Q',Q''$. Dazu wenden wir das ab Artikel 88 erläuterte Verfahren an und erhalten
\begin{center}
\begin{tabular}{@{}llll@{}}
aus & $r,\ r',\ v'-v,\ t'-t$ & den Wert & $(\eta01)$\\
& $r',\ r'',\ v''-v',\ t''-t'$ & . . . . . . . . & $(\eta12)$\\
& $r'',\ r''',\ v'''-v'',\ t'''-t''$ & . . . . . . . . & $(\eta23)$.
\end{tabular}
\end{center}
Diese Größen setzen wir zusammen mit den Werten von $r,r',r'',r''',\cos\tfrac12(v'-v)$ usw. in die Formeln III—VI ein. So ergeben sich für $P',Q',P'',Q''$ wesentlich genauere Werte als die, auf denen die erste Hypothese beruhte. Mit diesen Werten bildet man daher die zweite Hypothese. Führt man sie genau wie die erste zu Ende, so liefert sie für $P',Q',P'',Q''$ noch wesentlich genauere Werte und führt so zur dritten Hypothese. Diese Rechnungen wiederholt man, bis die Werte von $P',Q',P'',Q''$ keiner weiteren Korrektur mehr zu bedürfen scheinen. Dies richtig zu beurteilen lernt man bald durch häufige Übung. Ist die heliozentrische Bewegung gering, so liefert die erste Hypothese diese Werte meist schon mit hinreichender Genauigkeit. Umfasst die Bewegung dagegen einen größeren Bogen und weichen außerdem die Zeitabstände erheblich voneinander ab, so sind mehrere Hypothesen nacheinander durchzurechnen. In einem solchen Fall brauchen die Rechnungen für die ersten Hypothesen keine große Genauigkeit. In der letzten Hypothese bestimmt man schließlich die Elemente selbst auf die eben angegebene Weise.
